\documentclass[11pt,a4paper]{article}
\usepackage{inputenc}
\usepackage{fontenc}
\usepackage[ngerman]{babel}
\usepackage{amsmath,amssymb,amsthm}
\usepackage{mathtools}
\usepackage{geometry}
\geometry{margin=2.5cm}

\theoremstyle{definition}
\newtheorem{definition}{Definition}[section]
\newtheorem{proposition}{Proposition}[section]
\newtheorem{remark}{Bemerkung}[section]
\newtheorem{interpretation}{Interpretation}[section]

\title{Eine EABC-Kodierungstheorie für admissible Primzahl-Tupel}
\author{}
\date{13. April 2026}

\begin{document}
	\maketitle
	
	\begin{abstract}
		Admissible Primzahl-Tupel werden als Wörter über dem Alphabet $\mathcal{A}=\{E,A,B,C\}$ kodiert, das den Restklassen $1,5,7,11 \bmod 12$ entspricht. Arithmetische Zulässigkeit wirkt dabei als harter Filter auf den Symbolraum. Der Artikel entwickelt die symbolische Dynamik dieser Sprache, bestimmt ihre topologische und Shannon-Entropie und erweitert das Modell um geometrische Gewichtungen aus Kugel- und Randfiltern. Das Resultat ist eine Gibbs-Shannon-Struktur, in der Kodierlängen sowohl arithmetische als auch geometrische Selektivität abbilden.
	\end{abstract}
	
	\section{Motivation}
	
	Die bisherige Analyse admissibler Primzahl-Tupel legt nahe, dass diese nicht nur als arithmetische Konstellationen, sondern auch als \emph{strukturierte Symbolwörter} aufgefasst werden können. Grundlage ist die Zerlegung der Primzahlen \(p>3\) nach ihren Restklassen modulo \(12\):
	\[
	E \equiv 1 \pmod{12},\qquad
	A \equiv 5 \pmod{12},\qquad
	B \equiv 7 \pmod{12},\qquad
	C \equiv 11 \pmod{12}.
	\]
	Damit wird jedes Primzahl-Tupel zu einem Wort über dem Alphabet
	\[
	\mathcal A=\{E,A,B,C\}.
	\]
	
	Die entscheidende Beobachtung ist, dass admissible Primzahl-Tupel \emph{nicht} alle möglichen Wörter über \(\mathcal A\) realisieren, sondern nur eine stark eingeschränkte, hochstrukturierte Teilsprache. Diese Situation ist formal eng verwandt mit einer Shannon-artigen Kodierungstheorie: Eine Quelle produziert Zeichen, doch arithmetische Zulässigkeit wirkt als harter Filter auf den erlaubten Symbolraum.
	
	\section{Alphabet und Sprachraum}
	
	\begin{definition}[EABC-Alphabet]
		Es sei
		\[
		\mathcal A=\{E,A,B,C\}
		\]
		das Alphabet der vier zulässigen Restklassen modulo \(12\) für Primzahlen größer als \(3\), mit
		\[
		E \leftrightarrow 1,\qquad
		A \leftrightarrow 5,\qquad
		B \leftrightarrow 7,\qquad
		C \leftrightarrow 11 \pmod{12}.
		\]
	\end{definition}
	
	\begin{definition}[Kodierung eines Tupels]
		Für ein Primzahl-Tupel
		\[
		Q=(q_1,\dots,q_n)
		\]
		definieren wir seine EABC-Kodierung durch
		\[
		\Phi(Q)=(f(q_1),\dots,f(q_n))\in \mathcal A^n,
		\]
		wobei \(f(q_i)\in\{E,A,B,C\}\) die Familie der Primzahl \(q_i\) bezeichnet.
	\end{definition}
	
	\begin{definition}[Zulässige Wortmenge]
		Die Menge der durch admissible Primzahl-\(n\)-Tupel realisierten Wörter sei
		\[
		\mathcal W_n \subseteq \mathcal A^n.
		\]
		Die gesamte admissible Sprache ist dann
		\[
		\mathcal W=\bigcup_{n\ge 1}\mathcal W_n.
		\]
	\end{definition}
	
	Damit ist der Schritt von der Arithmetik zur Informationstheorie vollzogen: Primzahl-Tupel werden als Codewörter einer eingeschränkten symbolischen Quelle aufgefasst.
	
	\section{Beispiele und erste Struktur}
	
	Die minimale linke Kette dichter admissibler Tupel
	\[
	\{0,2\}
	\subset
	\{0,2,6\}
	\subset
	\{0,2,6,8\}
	\subset
	\{0,2,6,8,12\}
	\subset
	\{0,2,6,8,12,18\}
	\subset
	\{0,2,6,8,12,18,20\}
	\subset
	\{0,2,6,8,12,18,20,26\}
	\]
	liefert für geeignete Startklassen zyklische Familienwörter. Beispielsweise ergibt sich aus der Startklasse \(C\) die Folge
	\[
	CEABC,\quad CEABCA,\quad CEABCAB,\quad CEABCABE,\quad CEABCABEA,\dots
	\]
	Die zulässige Sprache ist somit weder zufällig noch frei, sondern besitzt eine eingebaute zyklische und geschachtelte Wortstruktur.
	
	\begin{remark}
		In einem freien Vierer-Alphabet gäbe es \(4^n\) mögliche Wörter der Länge \(n\). Die admissiblen Tupel realisieren dagegen nur einen winzigen Teil davon. Dies ist der erste Hinweis auf eine stark reduzierte Entropie.
	\end{remark}
	
	\section{Übergangsgraph und symbolische Dynamik}
	
	\begin{definition}[Übergangsgraph]
		Die admissible Sprache \(\mathcal W\) werde durch einen gerichteten Graphen
		\[
		G=(\mathcal A,\mathcal E)
		\]
		modelliert, wobei eine Kante
		\[
		X \to Y
		\]
		genau dann existiert, wenn der Übergang von Familie \(X\) zu Familie \(Y\) in einem admissiblen Tupelwort vorkommen kann.
	\end{definition}
	
	Im einfachsten zyklischen Modell gilt
	\[
	E\to A,\qquad A\to B,\qquad B\to C,\qquad C\to E.
	\]
	Dann erhält man die Adjazenzmatrix
	\[
	M=
	\begin{pmatrix}
		0&1&0&0\\
		0&0&1&0\\
		0&0&0&1\\
		1&0&0&0
	\end{pmatrix}
	\]
	in der Reihenfolge \((E,A,B,C)\).
	
	\begin{definition}[Zulässiges Wort]
		Ein Wort
		\[
		w=w_1w_2\cdots w_n \in \mathcal A^n
		\]
		heißt zulässig, wenn alle Übergänge
		\[
		w_i \to w_{i+1}
		\]
		im Graphen \(G\) erlaubt sind.
	\end{definition}
	
	Damit wird die EABC-Struktur zu einem Objekt der symbolischen Dynamik.
	
	\section{Symbolische Entropie}
	
	\begin{definition}[Shannon-Entropie auf \(\mathcal W_n\)]
		Sei \(p_n(w)\) eine Wahrscheinlichkeitsverteilung auf \(\mathcal W_n\). Dann definieren wir die symbolische Entropie
		\[
		H_n=-\sum_{w\in\mathcal W_n} p_n(w)\log p_n(w).
		\]
	\end{definition}
	
	\begin{definition}[Gleichverteilungsentropie]
		Falls alle zulässigen Wörter der Länge \(n\) gleich wahrscheinlich sind, also
		\[
		p_n(w)=\frac{1}{|\mathcal W_n|},
		\]
		dann gilt
		\[
		H_n=\log |\mathcal W_n|.
		\]
	\end{definition}
	
	\begin{definition}[Entropierate]
		Die asymptotische symbolische Entropierate sei
		\[
		h=\limsup_{n\to\infty}\frac{H_n}{n}.
		\]
	\end{definition}
	
	Im freien Fall über vier Symbolen wäre
	\[
	H_n^{\max}=n\log 4.
	\]
	Daraus ergibt sich als Maß der arithmetischen Einschränkung:
	
	\begin{definition}[Entropiereduktion]
		Die arithmetische Entropiereduktion ist
		\[
		\Delta H_n = n\log 4 - H_n.
		\]
	\end{definition}
	
	\begin{interpretation}
		\(\Delta H_n\) misst genau, wie stark die Zulässigkeitsbedingungen der Primzahl-Arithmetik die freie Symbolquelle einschränken.
	\end{interpretation}
	
	\section{Topologische Entropie}
	
	\begin{proposition}
		Ist die admissible Sprache durch eine endliche Adjazenzmatrix \(M\) beschrieben, dann ist ihre topologische Entropie
		\[
		h_{\mathrm{top}}=\log \lambda_{\max}(M),
		\]
		wobei \(\lambda_{\max}(M)\) der größte Eigenwert von \(M\) ist.
	\end{proposition}
	
	\begin{proof}[Skizze]
		Dies ist das Standardresultat aus der Theorie der Subshifts endlichen Typs: Die Anzahl der zulässigen Wörter wächst asymptotisch wie \(\lambda_{\max}(M)^n\). Durch Logarithmieren und Normieren auf \(n\) ergibt sich die Behauptung.
	\end{proof}
	
	\begin{remark}
		Im rein zyklischen Viererfall ist \(\lambda_{\max}(M)=1\), also
		\[
		h_{\mathrm{top}}=0.
		\]
		Das bedeutet: Die Sprache ist fast deterministisch, also maximal geordnet.
	\end{remark}
	
	\section{Geometrische Gewichtung der Wörter}
	
	Die symbolische Sprache allein erfasst nur die arithmetische Zulässigkeit. Die numerischen Experimente mit Kugel- und Randfiltern legen nahe, zusätzlich eine geometrische Gewichtung einzuführen.
	
	\begin{definition}[Vierlings-\(S^3\)-Term]
		Für einen Primzahlvierling \(Q=(q_1,q_2,q_3,q_4)\) definieren wir
		\[
		D_\Sigma = 4\,\frac{\sum_{i=1}^4 q_i^{3/4}}{\sum_{i=1}^4 q_i},
		\qquad
		S_4(Q)=\frac{1}{D_\Sigma}.
		\]
	\end{definition}
	
	\begin{definition}[Fünflings-Filter]
		Für einen Primzahlfünfling \(P=(p_1,\dots,p_5)\) definieren wir
		\[
		D_{5\to3}=\frac{15}{4}\,\frac{\sum_{i=1}^5 p_i^{3/5}}{\sum_{i=1}^5 p_i},
		\qquad
		S_5(P)=\frac{1}{D_{5\to3}}.
		\]
	\end{definition}
	
	\begin{definition}[Ternäre Korrektur]
		Für ein geeignet geordnetes Tupel definieren wir
		\[
		\chi
		=
		\log\!\left(
		\frac{(q_1q_2q_3)^{1/3}+(q_2q_3q_4)^{1/3}}{2}
		\right)
		-\log m,
		\]
		wobei \(m\) ein geeignetes Zentrum des Tupels ist.
	\end{definition}
	
	Diese Größen interpretieren wir als geometrische oder arithmetisch-geometrische \emph{Kosten} des zugehörigen Codeworts.
	
	\section{Gibbs-Shannon-Verteilung}
	
	\begin{definition}[Gewichtete Wortenergie]
		Jedem zulässigen Wort \(w\in\mathcal W_n\) werde eine Energie
		\[
		E(w)=E_{\mathrm{sym}}(w)+E_{\mathrm{geom}}(w)
		\]
		zugeordnet. Beispielsweise kann man setzen
		\[
		E_{\mathrm{geom}}(w)
		=
		\alpha\,\widetilde{S_4}
		+
		\beta\,\widetilde{S_5}
		+
		\gamma\,\widetilde{\chi},
		\]
		wobei die Tilden normierte Versionen der jeweiligen Größen bezeichnen.
	\end{definition}
	
	\begin{definition}[Gibbs-Shannon-Wahrscheinlichkeit]
		Die Wahrscheinlichkeit eines zulässigen Wortes \(w\in\mathcal W_n\) sei
		\[
		p_n(w)=\frac{e^{-\tau E(w)}}{Z_n},
		\qquad
		Z_n=\sum_{u\in\mathcal W_n} e^{-\tau E(u)},
		\]
		wobei \(\tau>0\) ein inverser Temperaturparameter ist.
	\end{definition}
	
	\begin{definition}[Arithmetisch-geometrische Entropie]
		Die daraus entstehende gewichtete Shannon-Entropie ist
		\[
		H_n^{\mathrm{arith\text{-}geom}}
		=
		-\sum_{w\in\mathcal W_n} p_n(w)\log p_n(w).
		\]
	\end{definition}
	
	\begin{interpretation}
		Diese Größe misst nicht nur die symbolische Freiheit der Quelle, sondern zugleich die geometrische Selektivität admissibler Tupel.
	\end{interpretation}
	
	\section{Kodierlängen}
	
	Nach Shannon ist die ideale Kodierlänge eines Wortes proportional zu seinem negativen Logarithmus der Wahrscheinlichkeit.
	
	\begin{definition}[Ideale Kodierlänge]
		Für ein zulässiges Wort \(w\) definieren wir
		\[
		L(w)\approx -\log p_n(w).
		\]
	\end{definition}
	
	Damit gilt:
	\begin{itemize}
		\item häufige und geometrisch billige Wörter sind kurz kodierbar,
		\item seltene oder geometrisch teure Wörter benötigen längere Codes.
	\end{itemize}
	
	Die Primzahl-Tupelsprache wird somit zu einer \emph{komprimierbaren, aber stark eingeschränkten Informationsquelle}.
	
	\section{Global-lokale Zerlegung}
	
	Die numerischen Befunde legen eine Zweiteilung nahe:
	
	\begin{itemize}
		\item \textbf{globale Struktur:} Fünflingsfilter wie \(S_5\) und ternäre Korrekturen \(\chi\),
		\item \textbf{lokale Struktur:} Vierlingsgrößen wie \(\delta\), \(\delta^3\) und \(S_4\).
	\end{itemize}
	
	Dies motiviert eine Zerlegung
	\[
	E(w)=E_{\mathrm{glob}}(w)+E_{\mathrm{lok}}(w),
	\]
	etwa mit
	\[
	E_{\mathrm{glob}}(w)
	=
	a_1\,\widetilde{S_5}
	+
	a_2\,\widetilde{\chi}
	+
	a_3\,\widetilde{(S_4\chi)},
	\]
	und
	\[
	E_{\mathrm{lok}}(w)
	=
	b_1\,\delta
	+
	b_2\,\delta^3
	+
	b_3\,\widetilde{S_4}.
	\]
	
	\begin{remark}
		Diese Zerlegung ist konzeptionell wichtig: Die bisherigen numerischen Tests deuten darauf hin, dass globale Spektralskalierung und lokale Abstandsstruktur nicht von derselben Größe getragen werden.
	\end{remark}
	
	\section{Zusammenfassung des Modells}
	
	Die EABC-Kodierungstheorie beruht damit auf vier Ebenen:
	\begin{enumerate}
		\item einem endlichen Alphabet \(\{E,A,B,C\}\),
		\item einer arithmetisch admissiblen Teilsprache \(\mathcal W\),
		\item einer symbolischen Entropie \(H_n\) bzw. topologischen Entropie \(h_{\mathrm{top}}\),
		\item einer geometrischen Gewichtung durch Kugel- und Randfilter.
	\end{enumerate}
	
	Der Kernsatz der Theorie lautet somit:
	
	\begin{quote}
		Die EABC-Sprache admissibler Primzahl-Tupel ist eine symbolische Quelle mit arithmetisch reduzierter Entropie und geometrisch gewichteter Kodierstruktur.
	\end{quote}
	
	\section{Ausblick}
	
	Die nächsten mathematisch natürlichen Schritte sind:
	\begin{itemize}
		\item explizite Bestimmung der Übergangsgraphen für \(4\)-, \(5\)-, \(6\)-, \(7\)-, \(8\)- und \(9\)-Tupel,
		\item Berechnung der zugehörigen topologischen Entropien,
		\item Aufbau einer gewichteten Gibbs-Struktur mit \(S_4\), \(S_5\) und \(\chi\),
		\item Vergleich der resultierenden Kodierlängen mit den bereits beobachteten spektralen Korrelationsmustern.
	\end{itemize}
	
	Damit entsteht eine erste Brücke zwischen
	\begin{center}
		\emph{admissibler Primzahl-Arithmetik, symbolischer Dynamik, Entropietheorie und geometrischer Filterstruktur.}
	\end{center}
	
\end{document}